\documentclass[12pt]{article}
\usepackage[utf8]{inputenc}
\usepackage[T1]{fontenc}
\usepackage[margin=1in]{geometry}
\usepackage{setspace}
\usepackage{xcolor}
\usepackage[colorlinks=true, urlcolor=blue!60!black, linkcolor=blue!60!black]{hyperref}
\setlength{\parindent}{2em}
\setlength{\parskip}{1.25em}
\renewcommand{\baselinestretch}{1.0}
\begin{document}
\begin{center}
{\LARGE \textit{South Dakota v. Dole}, 483 U.S. 203 (1987)\par}
\vspace{0.8em}
{\large Case Brief --- Four-Page Edition\par}
\end{center}
\vspace{1em}

\section{Memory Jogger}

Upholding 7--2 a federal law withholding 5\% of specified highway funds from States that allowed purchase or public possession of alcohol by people under 21, the Court set the test for conditional spending: conditions must serve the general welfare, be unambiguous, relate to the federal interest, violate no independent constitutional bar, and not coerce (\href{https://supreme.justia.com/cases/federal/us/483/203/}{483 U.S. 203}).

\section{Detailed Case Facts}

In 1984 Congress enacted 23 U.S.C. \S~158, the National Minimum Drinking Age Act, after the Presidential Commission on Drunk Driving found that inconsistent state drinking ages created ``an incentive to drink and drive'' by encouraging young adults to cross state lines to buy alcohol. Rather than banning underage sales directly, \S~158 directed the Secretary of Transportation to withhold 5\% of certain federal highway funds from any State in which purchase or public possession of alcohol by a person under 21 was lawful. South Dakota allowed 19-year-olds to buy beer containing up to 3.2\% alcohol.

\section{Procedural History}

South Dakota sued Secretary of Transportation Elizabeth Dole, claiming that \S~158 exceeded the Spending Clause and violated the Twenty-first Amendment's reservation of state power over alcohol. The District Court rejected the claims and the Eighth Circuit affirmed. The Supreme Court granted certiorari, heard argument on April 28, 1987, and affirmed on June 23, 1987.

\section{Judicial Votes}

Affirmed, \textbf{7--2}.
\begin{itemize}
\item \textbf{Majority:} Chief Justice Rehnquist, joined by Justices White, Marshall, Blackmun, Powell, Stevens, and Scalia.
\item \textbf{Concurrences:} None.
\item \textbf{Dissents:} Justice Brennan; Justice O'Connor.
\end{itemize}

\section{Holding}

Even assuming the Twenty-first Amendment bars Congress from setting a national drinking age directly (a question the Court did not decide), \S~158 is a valid exercise of the spending power. The condition (1) pursues the general welfare; (2) is stated unambiguously; (3) is directly related to the federal interest in safe interstate travel; (4) violates no ``independent constitutional bar,'' because it does not induce States to act unconstitutionally; and (5) is ``relatively mild encouragement,'' not compulsion.

\section{Analysis of Opinions}

\textbf{Majority (Rehnquist).} Sidestepping the Twenty-first Amendment question, Rehnquist relied on \textit{United States v. Butler}'s Hamiltonian view that the spending power ``is not limited by the direct grants of legislative power found in the Constitution.'' Courts defer substantially to Congress on the general welfare, and the Court doubted it is judicially enforceable at all. Under \textit{Pennhurst}, conditions must let States choose ``knowingly,'' and \S~158 was clear. Without defining germaneness's limits, the Court found the condition ``directly related to one of the main purposes for which highway funds are expended---safe interstate travel.'' The independent-bar limit does not forbid indirect achievement of objectives beyond Congress's direct powers; it forbids inducing unconstitutional conduct such as invidious discrimination, and a State could lawfully raise its drinking age. Finally, losing 5\% of specified highway funds (later described in \textit{NFIB} as less than half of one percent of South Dakota's budget) was not coercive, and success does not prove compulsion.

\textbf{Dissent (Brennan).} Setting the drinking age is reserved to the States by the Twenty-first Amendment, so Congress cannot condition a grant in a way that abridges that power.

\textbf{Dissent (O'Connor).} Accepting the framework, O'Connor dissented on relatedness. The condition was ``far too over- and under-inclusive'': it barred teenagers from drinking even when not driving, while 84\% of alcohol-related fatalities involved drivers over 21. Congress may specify how federal money is spent but may not regulate beyond the expenditure; \S~158 was a regulation of who may drink, and the majority's approach would let Congress reach almost any aspect of state life.

\textbf{Later developments.} \textit{NFIB v. Sebelius} (2012) found the ACA's threat to all existing Medicaid funds ``a gun to the head''; \textit{Landor} (2026) held that \textit{Dole}'s requirements apply ``in addition to---not instead of'' the rule that Spending Clause statutes bind only those who knowingly consent.

\section{Examples: Future Applications}

\begin{enumerate}
\item \textbf{Same-Side (Valid): New EV-Corridor Grants.} A new grant program requiring participants to build interstate chargers and adopt EV-ready building codes is related, clear, and noncoercive, since nonparticipants lose no existing funds.
\item \textbf{Same-Side (Valid): Blood-Alcohol Limit.} Withholding 5\% of highway funds from States without a 0.08 limit is even more closely related to highway safety than the drinking-age condition.
\item \textbf{Opposite-Side (Invalid): All K--12 Funds for a Campus Gun Ban.} Threatening 100\% of existing Title I funds, 8\% of a typical state budget, is ``a gun to the head'' under \textit{NFIB}, and the link to university gun policy is doubtful.
\item \textbf{Opposite-Side (Invalid): Racial Admissions Quotas.} A small, optional grant conditioned on racial quotas fails the independent-bar prong, because a State could not adopt a quota itself.
\item \textbf{Fence-Sitter: 25\% of Education Funds for a Phone Ban.} The condition is clear and related, but the threatened share falls between \textit{Dole}'s 5\% and \textit{NFIB}'s full Medicaid stake, and neither case supplies a metric.
\end{enumerate}

\section{Critique}

\textbf{A ``Trojan Horse'' for enumerated powers.} Scholars note the irony that Rehnquist, architect of the \textit{Lopez} and \textit{Morrison} federalism revival, wrote the opinion letting Congress buy state compliance with policies it could not mandate, while federal taxes on the same taxpayers press States to accept the money.

\textbf{Indeterminacy and asymmetry.} Neither \textit{Dole} nor \textit{NFIB} supplies a bright line for relatedness or coercion, and footnote 3 left germaneness undefined. The independent-bar test protects individual rights from indirect subversion but treats structural federalism protections as waivable for cash.

\section{Key Quotations}

\begin{quote}
``These cases establish that the `independent constitutional bar' limitation on the spending power is not, as petitioner suggests, a prohibition on the indirect achievement of objectives which Congress is not empowered to achieve directly. Instead, we think that the language in our earlier opinions stands for the unexceptionable proposition that the power may not be used to induce the States to engage in activities that would themselves be unconstitutional.''
\end{quote}
\noindent --- Chief Justice Rehnquist, Opinion of the Court, 483 U.S. at 210.

\begin{quote}
``Thus, objectives not thought to be within Article I's `enumerated legislative fields,' \textit{id.} at 297 U. S. 65, may nevertheless be attained through the use of the spending power and the conditional grant of federal funds.''
\end{quote}
\noindent --- Chief Justice Rehnquist, Opinion of the Court, 483 U.S. at 207.

\begin{quote}
``Our decisions have recognized that, in some circumstances, the financial inducement offered by Congress might be so coercive as to pass the point at which `pressure turns into compulsion.' . . . Here Congress has offered relatively mild encouragement to the States to enact higher minimum drinking ages than they would otherwise choose. But the enactment of such laws remains the prerogative of the States not merely in theory, but in fact.''
\end{quote}
\noindent --- Chief Justice Rehnquist, Opinion of the Court, 483 U.S. at 211--212.

\end{document}
